\section{Continuous-law qualification in the investment model}\label{sec:continuous95}
The risk theorem becomes a numerical method only when its inequalities can be constructed from economic primitives. We now supply this construction for the original investment technology with a nested tail-risk preference. States and innovations remain continuous, every feasible action is covered, and the implemented terminal controller rather than a hypothetical exact optimizer determines policy cost. This application joins continuation construction, the original-optimum comparison and a representation comparison on one problem. It does not replace the general NBO construction or the other economic applications.

\subsection{Primitives and the object being approximated}
Let $d\geq2$ be even, $X=[0,1]^d$, and interpret subscripts cyclically. At dates $t=0,1$, the feasible investment is $0\leq a\leq A(x)=1/8+\bar x/8$, where $\bar x=d^{-1}\sum_jx_j$. Write $b_j=1/2$ for odd $j$ and $b_j=1/4$ for even $j$, and $s_j=(-1)^{j+1}$. The original transition is
\begin{equation}\label{eq:investment95}
 F_j(x,a,z)=\frac1{16}+\frac{x_j}{2}+\frac{x_{j+1}}8
       +\frac{x_j(1-x_{j+1})}{16}+b_ja+s_jz,
 \qquad z\sim\operatorname{Unif}[-1/32,1/32].
\end{equation}
Define
\begin{align}\label{eq:cost95}
 S(x)&=\frac2d\sum_j(x_j-5/8)^2+
       \frac1{4d}\sum_j(x_j-x_{j+1})^2+2(1/2-2\bar x)_+^2,\\
 g(x)&=S(x)+\frac2d\sum_j(x_j-5/8)^2,\qquad h(a)=a^2+4a^4.\nonumber
\end{align}
We use $\beta=15/16$ and, separately, $\alpha\in\{1/2,3/4\}$. Average value at risk is the upper-tail cost functional
\[
 \rho_\alpha(v)=\sup\left\{\int vq\,dP:
             0\leq q\leq(1-\alpha)^{-1},\ \int q\,dP=1\right\}.
\]
Here $P$ is the original uniform innovation law. The objective is recursively defined by $V_2=g$ and
\[
 V_t(x)=S(x)+\inf_{0\leq a\leq A(x)}
       \{h(a)+\beta\rho_\alpha(V_{t+1}(F(x,a,\cdot)))\}.
\]
It is not the static tail risk of total realized cost. A complete policy value replaces each infimum by the action prescribed by that policy, with the same recursive risk primitive.

\paragraph{Assumptions for the continuous construction.}
For this subsection, the primitives are exactly \eqref{eq:investment95}--\eqref{eq:cost95}; $\alpha$ is fixed before construction; the fitted continuation $f$ is a finite convex function on $X$; and every numerical lower or upper value is an outward enclosure. Feasible tail densities and feasible actions are retained explicitly. Terminal selection is the measurable, dyadically acquired controller described below, with its fallback included in the error bound. These concrete conditions imply compact feasible graphs, bounded costs, measurable minimizers, and the conditional-risk dual used in Theorem~\ref{thm:centeredrisk92}. They do not assert that every recursive preference or every trained network has this convex structure.

\begin{proposition}[Primitive curvature and state supports]\label{prop:curvature95}
For either risk level and every even $d$, the one-date objective
$Q_1(x,a)=S(x)+h(a)+\beta\rho_\alpha(g(F(x,a,\cdot)))$
is jointly convex on $X\times[0,1/4]$, with curvature bounded below by
\begin{equation}\label{eq:curvature95}
 \frac{4-5\beta/4}{d}\|v\|_2^2+2w^2
\end{equation}
in a direction $(v,w)$. Its value $V_1$ is convex on $X$ and satisfies
\begin{equation}\label{eq:lipschitz95}
 |V_1(x)-V_1(y)|\leq\frac{15}{d}\|x-y\|_1.
\end{equation}
Let $W$ be the objective evaluated under any fixed feasible risk density, at an anchor $(x_0,a_0)$, and let $(p_x,p_a)$ be its gradient there. For any $\lambda,\nu\geq0$, set $e=p_a+\lambda-\nu$. Then the following is a global affine minorant of $V_1$:
\begin{align}\label{eq:affine95}
 \ell(y)={}&W+\lambda(a_0-A(x_0))-\nu a_0-\frac{e^2}{4}\nonumber\\
 &+\left(p_x-\frac{\lambda}{8d}\boldsymbol{1}\right)^{\!\top}(y-x_0).
\end{align}
Neither optimal multipliers nor an exact optimizing action are required.
\end{proposition}

The curvature is obtained from the original nonlinear technology, not imposed on an interpolated value grid. The negative second-order term from $x_j(1-x_{j+1})$ is bounded using $|\partial_jg|\leq10/d$; the stage curvature dominates it. The allowance $e^2/4$ pays for imperfect stationarity. The proof in the Technical Supplement also establishes the global state modulus in \eqref{eq:lipschitz95}, including the state-dependent action constraint.

\subsection{An all-action certificate under the original innovation law}
Partition the innovation interval into $n$ equal bins $B_i$, with midpoints $z_i$ and probabilities $p_i=1/n$. An innovation bin is a set over which an inequality holds, not a replacement point mass in the economic model. Choose nonnegative $q_i\leq p_i/(1-\alpha)$ summing to one. These masses define a feasible density constant on each bin. For a fixed initial state and feasible proposal $a_0$, let $W_0$ be a lower bound on the objective under this density, obtained from globally valid affine supports of the continuation at the successor midpoints. Let $p$ enclose the derivative of the same supporting objective at $a_0$. Then
\begin{equation}\label{eq:actionlower95}
 \underline V_0(x)=\inf_{0\leq u\leq A(x)}
         \{W_0+p(u-a_0)+(u-a_0)^2\}
\end{equation}
is a lower bound on the original Bellman optimum when the supported continuation is $V_1$. The scalar quadratic is minimized over the full original interval. Interval uncertainty in $p$ is covered by taking the smaller minimum at its two endpoints, since the minimum is concave in $p$.

For a convex fitted continuation, midpoint supports give the lower risk witness and the larger endpoint value bounds the continuation throughout each bin. For $V_1$, Proposition~\ref{prop:curvature95} supplies midpoint supports and the allowance $15/(32n)$ covers each half-bin. Upper tail risk is the average of the largest $n(1-\alpha)$ upper bin bounds. This average is the exact risk of a dominating step function, hence an upper bound under the continuous law.

\begin{theorem}[Constructed original-optimum qualification]\label{thm:continuous95}
Suppose the supports and bin bounds above are outwardly verified. Suppose further that the complete terminal controller satisfies
$V_1\leq V_1^{\widehat\pi}\leq V_1+\varepsilon_1$ everywhere on $X$.
For every queried initial state, the lower value \eqref{eq:actionlower95} and the binwise upper bound $\overline J_0^{\widehat\pi}$ satisfy
\begin{equation}\label{eq:originalgap95}
 0\leq J_0^{\widehat\pi}(x)-V_0(x)
       \leq\overline J_0^{\widehat\pi}(x)-\underline V_0(x).
\end{equation}
For any stored convex continuation $f$, the same lower and upper constructions are admissible bounds for the centered-risk comparison of Theorem~\ref{thm:centeredrisk92}. Their validity does not depend on the representation of $f$, on a successful fit, or on the proposal being an exact minimizer.
\end{theorem}

To see the connection, write a verified continuation bracket as $\ell\leq V_1\leq u$ and use errors $L=\ell-f$, $U=u-f$. The centered identity gives
\[
 \rho_\alpha(f)+\mathcal E^f(L)=\rho_\alpha(\ell),\qquad
 \rho_\alpha(f)+\mathcal E^f(U)=\rho_\alpha(u).
\]
Joint evaluation avoids subtracting and then adding two separately rounded copies of $\rho_\alpha(f)$. The bound \eqref{eq:actionlower95} uses one feasible risk density to minorize the required all-action infimum; the policy upper bound controls the full risk envelope. Thus the application constructs sufficient inequalities for the original theorem. It does not claim that a non-neural continuation would receive a different validity theorem, or that computing the original-value witness is free.

The terminal controller observes each state coordinate rounded down to $40$ binary places. It accepts its floating proposal only after an outward all-action gap of at most $2^{-21}$ at that observed state; otherwise it uses rational bisection. The original-state transfer is less than $30\,2^{-40}$. Consequently $\varepsilon_1=2^{-20}$ covers the complete controller, including observation and fallback. Policy evaluation encloses this controller at every possible successor, not only at the midpoint states used by the integrator.

